\section{Discussion and limitations}
\label{sec:limits}

\paragraph{Setting and claims.}
The pipeline is an offline recognition stage: it assumes that the location of every instrument is given, and it is meant for reviewing archives, assisting annotation and quality control.
We make no claim of real-time or intraoperative use (at the 29\% operating point an instrument takes 3.0\,s on average on one A100), of state-of-the-art results, of label efficiency, or of robustness beyond the transfer test of Section~\ref{sec:shift}.
Published results are given for context only; they find their own boxes and ours are given, so the two are not the same task.

\paragraph{Evidence.}
Everything but the transfer test is measured on one dataset with five test cases, so intervals, mcIoU in particular, are wide, and each segmenter was trained once.
The uncertainty comparison is on a held-out fold of four cases; on the test cases only the evidential score is scored, and it finds fewer errors there (84.8\% at 20\% review) than on the fold (90.2\%).
\todo{update with the segmenter seed variance and the intervals of the held-out-threshold row}

\paragraph{What the multi-frame step does.}
At one frame per second the propagated mask sometimes leaves the instrument, yet a static view of the same region lowers the scores below the single pass, so the gain needs the tracker; the control does not separate following one instrument from combining views of the scene.
A closed Laparoscopic Grasper and a Suction Instrument look alike in a single frame and refinement does not separate them.
\todo{confirm that the count of the look-alike pairs (59 of the 210 errors of the earlier vote ensemble) still describes the final ensemble, or recompute it}

\paragraph{Distribution shift.}
The score degrades under zero-shot transfer to EndoVis, so a review budget set on GraSP would miss most errors on a new dataset; detecting shift is a separate problem that this score does not solve.
\todo{explore causal mode: a variant that reads only past frames}

\section{Conclusion}
\label{sec:conclusion}
Checking one instrument in five finds most of the errors of an instrument recogniser, and following the same kind of instruments through neighbouring frames recovers nearly all of the gain we measured at a fraction of the cost.
One pass of a four-network evidential ensemble is enough to rank them, matching a 12-network deep ensemble and beating 80-pass MC dropout at a fixed budget, though the routing degrades under transfer.
\todo{one sentence on what this enables, and on future work: a detector in front}
